\documentclass[10pt,a4paper]{amsart}
\usepackage[margin=19mm]{geometry}
\usepackage{amsmath,amssymb,mathtools}
\usepackage[scr=rsfs,cal=euler]{mathalfa}
\usepackage{xfrac}
\usepackage[unicode,hidelinks,bookmarks=false]{hyperref}
\newcommand{\ie}{i.e.\ }
\newcommand{\cf}{cf.\ }
\newcommand{\resp}{respectively\ }
\newcommand{\Subsection}{\subsection}
\providecommand{\nobreakdash}{\nobreak\hbox{-}}
\setlength{\emergencystretch}{2em}
\multlinegap0pt
\usepackage{bm}
\usepackage{bbm}
\usepackage{dsfont}
\usepackage{xspace}
\usepackage{cancel}
\usepackage{mathtools}
\usepackage{cleveref}
\usepackage{amsthm}
\makeatletter
\@ifundefined{theorem}{\newtheorem{theorem}{Theorem}}{}
\@ifundefined{lemma}{\newtheorem{lemma}[theorem]{Lemma}}{}
\@ifundefined{proposition}{\newtheorem{proposition}[theorem]{Proposition}}{}
\makeatother
\newcommand{\m}{\mathfrak m}
\newcommand{\q}{\mathfrak q}
\title[Dimensions of a ring and its formal power series ring]{Dimensions of a ring and its formal power series ring}
\author{Viet-Hoang Tran, Thieu N. Vo, Tan M. Nguyen}
\date{Source: arXiv:2609.17854v1 (CC BY 4.0)}
\begin{document}
\maketitle

\noindent\emph{Source and licence.} Dimensions of a ring and its formal power series ring. Version arXiv:2609.17854v1 (CC BY 4.0), distributed under the Creative Commons Attribution 4.0 International licence, as recorded in the arXiv licence metadata for this exact version and shown on the abstract page https://arxiv.org/abs/2609.17854v1. Full rights notice: licenses/arxiv-2609.17854-CC-BY-4.0.txt.\par\smallskip

\noindent\textit{Editorial scope.} Counted unit: the Proposition whose source label is \texttt{prop:lower} in the article (source0.tex lines 610--613, source label \texttt{prop:lower}), delivered together with the article's own complete original proof of it (source0.tex lines 615--671, one \texttt{proof} environment of 57 lines). No mathematics is written, repaired or re-derived: statement and proof are the article's own text line for line. Required context retained uncounted: eq:array (lines 540--543); prop:upper (lines 551--553); lem:exponentials (lines 578--581). Every definition, display and label of those lines is retained unchanged. Omitted from this package and not redistributed: 1--539, 544--550, 554--577, 582--609, 614--614, 672--759. Nothing inside a retained excerpt is omitted or altered: every retained body is delivered exactly as the article's lines are written, so no omission receipt of any kind accompanies this package. Citations: the retained excerpts carry no citation command at all, so no bibliography entry of the article is transcribed and no reference is redistributed. Reference adaptation: none was needed and none was made; every reference inside the delivered unit resolves inside it, so no unresolved reference or citation is produced. Printed slips kept literal: no printed word, symbol, hypothesis or number of the counted statement or of its proof was altered. Layout and class: the article's own class is \texttt{amsart} with packages including fontenc, lmodern, microtype, needspace, geometry, amsmath, amssymb, mathtools, amsthm, xcolor, hyperref, cleveref; the wrapper is the lane's pinned \texttt{amsart} editorial wrapper at 10pt on A4, which restates the theorem-like environments the retained text needs and adds the article's own macro definitions that the retained text needs (\texttt{\textbackslash m}, \texttt{\textbackslash q}), with extra wrapper packages (bm, bbm, dsfont, xspace, cancel, mathtools, cleveref). The delivered numbering is the unit's own. Assets: no retained span contains a figure environment, a graphics-inclusion command, a PDF-page inclusion, a table or a TikZ picture; no image file is copied into this package, the package declares no asset dependency and no image file is redistributed with it. Licence: version arXiv:2609.17854v1 of this article is distributed under the Creative Commons Attribution 4.0 International licence, as recorded in the arXiv licence metadata for this exact version and shown on the abstract page https://arxiv.org/abs/2609.17854v1. A full rights notice is delivered at licenses/arxiv-2609.17854-CC-BY-4.0.txt. Characters: every retained excerpt is pure ASCII, so the delivered engine, XeLaTeX with its default Latin Modern text font, reports no missing character.
\begin{equation}\label{eq:array}
 A=\left\{\sum_{n\ge0}b_nc^n\in B[[c]]:
                         b_n\in E_{2^n-1}\text{ for all }n\right\}.
\end{equation}

\begin{proposition}\label{prop:upper}
One has $\dim A\le r+2$.
\end{proposition}

\begin{lemma}\label{lem:exponentials}
There exist $f_1,\ldots,f_r\in xk[[x]]$ algebraically
independent over $k$.
\end{lemma}

\begin{proposition}\label{prop:lower}
The ring $A$ contains a prime chain of length $r+2$.
Thus $\dim T[[x]]=r+2$.
\end{proposition}

\begin{proof}
Choose the series in \cref{lem:exponentials} and put
\[
 \q_j=(u_1-f_1(x),\ldots,u_j-f_j(x))B\quad(1\le j\le r),
 \qquad\q_0=(0).
\]
Before localizing, substitution of $f_i(x)$ for $u_i$, $i\le j$,
defines a surjection
\[
 F[U]\longrightarrow F[u_{j+1},\ldots,u_r]
\]
with kernel $(u_1-f_1(x),\ldots,u_j-f_j(x))F[U]$.
Its kernel is disjoint from $S$: expanding a nonzero element of
$k[U]$ in the remaining variables would otherwise give a nonzero
polynomial relation over $k$ among $f_1,\ldots,f_j$.
Localization therefore gives proper prime ideals and a strict chain
\begin{equation}\label{eq:qchain}
 (0)=\q_0\subsetneq\q_1\subsetneq\cdots\subsetneq\q_r.
\end{equation}
For strictness, $u_j-f_j(x)$ is nonzero in the domain obtained
after the preceding substitution, and remains nonzero after
localization. Full substitution extends to $B\longrightarrow L$
with kernel $\q_r$, since every element of $S$ has nonzero image
in $F\subseteq L$. This map fixes $F$, so
\begin{equation}\label{eq:qF}
 \q_r\cap F=(0).
\end{equation}

The coefficientwise ideal $\q_j[[c]]\subseteq B[[c]]$ is prime,
since its quotient is $(B/\q_j)[[c]]$, a domain.
Contract these primes to $A$:
\[
 Q_j=\q_j[[c]]\cap A\qquad(0\le j\le r).
\]
They all avoid $c$, and $Q_0=(0)$.
Since $u_j=D_1u_j/D_1\in W_1$ and $f_j(x)\in F$, one has
\[
 c(u_j-f_j(x))\in Q_j\setminus Q_{j-1}
 \qquad(1\le j\le r),
\]
where membership in $A$ follows from~\eqref{eq:array}.
Thus the contracted chain is strict.

If $a\in Q_r$, its coefficient at $c^0$ belongs to $E_0=F$
by~\eqref{eq:array} and to $\q_r$ by definition. It is zero
by~\eqref{eq:qF}. This says exactly that every $x$-coefficient
of $a$ belongs to $\m$, so $Q_r\subseteq P=\m[[x]]$.
The inclusion is strict because $c\in P\setminus Q_r$.
Finally $P\subsetneq P+(x)$ is a strict prime inclusion, since
$A/P\cong k[[x]]$. Hence
\begin{equation}\label{eq:final-chain}
 (0)=Q_0\subsetneq Q_1\subsetneq\cdots\subsetneq Q_r
       \subsetneq P\subsetneq P+(x)
\end{equation}
has length $r+2$. Combining this with \cref{prop:upper} proves
the dimension equality.
\end{proof}

\end{document}
